\section{Three-city SEIR network}
\label{sec:network}

\subsection{Transmission support and orientation}

Let
\[
\mathcalV=\{1,2,3\}
\]
represent three cities. Cities 1 and 3 communicate only through city 2, and
each city also contains a within-city transmission channel. The effective
transmission support is ordered as
\[
\mathcalE=
\bigl(
(1,1),(1,2),(2,1),(2,2),(2,3),(3,2),(3,3)
\bigr).
\]
Under the convention used throughout the project, the ordered pair
$(i,j)\in\mathcalE$ corresponds to the matrix entry $a_{ij}$ and represents
epidemiological influence from source city $j$ to receiving city $i$.
Consequently, the graph-theoretic arrow points from $j$ to $i$.

A diagonal channel $(i,i)$ does not describe an individual infecting itself.
It represents aggregate contact between susceptible and infectious individuals
belonging to the same city.

\begin{figure}[ht]
\centering
\begin{tikzpicture}[
  city/.style={circle,draw=deepblue,fill=softblue,minimum size=13mm,font=\bfseries},
  inter/.style={-{Stealth[length=2.2mm]},line width=0.75pt,deepblue},
  intra/.style={-{Stealth[length=2mm]},line width=0.75pt,deepblue},
  lab/.style={font=\small,align=center}
]
\node[city] (c1) at (0,0) {1};
\node[city] (c2) at (3.3,0) {2};
\node[city] (c3) at (6.6,0) {3};

\draw[inter,bend left=18] (c1) to node[above,lab] {$a_{21}$} (c2);
\draw[inter,bend left=18] (c2) to node[below,lab] {$a_{12}$} (c1);
\draw[inter,bend left=18] (c2) to node[above,lab] {$a_{32}$} (c3);
\draw[inter,bend left=18] (c3) to node[below,lab] {$a_{23}$} (c2);

\draw[intra] (c1) edge[loop above,min distance=13mm] node[above,lab] {$a_{11}$} (c1);
\draw[intra] (c2) edge[loop above,min distance=13mm] node[above,lab] {$a_{22}$} (c2);
\draw[intra] (c3) edge[loop above,min distance=13mm] node[above,lab] {$a_{33}$} (c3);

\node[lab,below=12mm of c2] {Arrows show epidemiological influence from source to receiver.};
\end{tikzpicture}
\caption{Three-city mobility corridor with one diagonal transmission channel in each city.}
\label{fig:three_city_network}
\end{figure}

The incoming source sets are
\[
\Min(1)=\{1,2\},
\qquad
\Min(2)=\{1,2,3\},
\qquad
\Min(3)=\{2,3\}.
\]

\subsection{State and differential equations}

For each city $i$, let $S_i,E_i,I_i,R_i$ denote the susceptible, exposed,
infectious, and removed fractions. The state is ordered as
\[
x=\col
(S_1,S_2,S_3,E_1,E_2,E_3,I_1,I_2,I_3,R_1,R_2,R_3)
\in\R^{12}.
\]
The city-specific equations are
\[
\begin{aligned}
\dot S_1
&=-\beta S_1(a_{11}I_1+a_{12}I_2),
&
\dot E_1
&=\beta S_1(a_{11}I_1+a_{12}I_2)-\sigma E_1,\\
\dot S_2
&=-\beta S_2(a_{21}I_1+a_{22}I_2+a_{23}I_3),
&
\dot E_2
&=\beta S_2(a_{21}I_1+a_{22}I_2+a_{23}I_3)-\sigma E_2,\\
\dot S_3
&=-\beta S_3(a_{32}I_2+a_{33}I_3),
&
\dot E_3
&=\beta S_3(a_{32}I_2+a_{33}I_3)-\sigma E_3,
\end{aligned}
\]
with
\[
\dot I_i=\sigma E_i-\gamma I_i,
\qquad
\dot R_i=\gamma I_i,
\qquad i=1,2,3.
\]

\subsection{Exact quadratic representation}

Define the seven infection products in the fixed channel order:
\[
\phi^{SI}(x)=\col(
S_1I_1,S_1I_2,
S_2I_1,S_2I_2,S_2I_3,
S_3I_2,S_3I_3).
\]
Let
\[
B_0=
\begin{bmatrix}
 a_{11} & a_{12} & 0 & 0 & 0 & 0 & 0\\
 0 & 0 & a_{21} & a_{22} & a_{23} & 0 & 0\\
 0 & 0 & 0 & 0 & 0 & a_{32} & a_{33}
\end{bmatrix}.
\]
Then the quadratic transmission contribution is represented by
\[
f_{\mathrm Q}(x)
=
\beta
\begin{bmatrix}
-B_0\\
 B_0\\
 0_{3\times7}\\
 0_{3\times7}
\end{bmatrix}
\phi^{SI}(x).
\]
The remaining compartmental transitions form a linear field
\[
f_{\mathrm L}(x)=A_{\ell}x.
\]
Thus the complete vector field is decomposed as
\[
f(x)=f_{\mathrm L}(x)+f_{\mathrm Q}(x).
\]
This decomposition is not merely notational. It separates the two operations
used by Algorithm~1:

\begin{itemize}[leftmargin=2em]
\item $f_{\mathrm Q}$ increases polynomial degree by one;
\item $f_{\mathrm L}$ preserves degree and therefore determines the linear
      closure required within each homogeneous block.
\end{itemize}
